\documentclass[11pt]{article}
\usepackage{assets/tex/notes/eric-notes-style}

\newcommand{\R}{\mathbb R}
\newcommand{\N}{\mathbb N}
\newcommand{\norm}[1]{\left\lVert #1\right\rVert}

\begin{document}

\notetitle{MATH 320-3: Real Analysis III Notes}{Topology, Differentiability, and Measure}

\topic{Compactness}

\entry{Covers}
\[
\text{Let }E\subseteq X.
\]
\[
\text{An open cover of }E\text{ is }\mathcal V=\{V_\alpha:\alpha\in A\}
\text{ of open sets }V_\alpha\text{ in }X
\]
\[
\text{such that }E\subseteq \bigcup_{\alpha\in A}V_\alpha.
\]
\[
\mathcal V\text{ has a finite subcover if }\exists\text{ a finite set }A_0\subseteq A
\text{ such that }E\subseteq\bigcup_{\alpha\in A_0}V_\alpha.
\]

\entry{Compactness}
\[
E\text{ is compact}\Longleftrightarrow
\text{every open cover of }E\text{ has a finite subcover.}
\]

\entry{Compact iff sequentially compact}
\[
\text{In a metric space }X,\ E\subseteq X,\qquad
E\text{ is compact}\Longleftrightarrow E\text{ is sequentially compact.}
\]

\entry{Heine-Borel}
\[
\text{In a metric space }X\text{ with the Bolzano-Weierstrass property, }E\subseteq X,
\]
\[
E\text{ is compact}\Longleftrightarrow E\text{ is closed and bounded.}
\]

\entry{Uniform Continuity}
\[
f\text{ continuous on compact }E
\Longrightarrow f\text{ is uniformly continuous on }E.
\]

\entry{Continuity preserves compactness}
\[
f\text{ is continuous on compact }E
\Longrightarrow f(E)\text{ is compact in }Y.
\]

\entry{Continuity and pre-image}
\[
f\text{ is continuous on }E
\Longleftrightarrow f^{-1}(U)\text{ is open in }E,\quad \forall U\text{ open in }Y,
\]
\[
\Longleftrightarrow f^{-1}(B)\text{ is closed in }E,\quad \forall B\text{ closed in }Y.
\]

\entry{Extreme Value Theorem}
\[
\varnothing\ne E\text{ compact},\qquad f:E\to\R\text{ continuous}
\]
\[
\Longrightarrow
\exists x_{\min},x_{\max}\in E
\text{ such that }f(x_{\min})=\min_E f,\quad f(x_{\max})=\max_E f.
\]

\entry{Continuity of inverse function}
\[
f\text{ continuous on compact }E,\ f\text{ one-to-one}
\Longrightarrow f^{-1}\text{ is continuous.}
\]

\entry{Images and pre-images}
\[
\text{Let }X,Y\text{ be sets and }f:X\to Y.
\]
\[
f\left(\bigcup_{\alpha\in A}E_\alpha\right)=\bigcup_{\alpha\in A}f(E_\alpha),
\qquad
f\left(\bigcap_{\alpha\in A}E_\alpha\right)\subseteq\bigcap_{\alpha\in A}f(E_\alpha).
\]
\[
f^{-1}\left(\bigcup_{\alpha\in A}B_\alpha\right)=\bigcup_{\alpha\in A}f^{-1}(B_\alpha),
\qquad
f^{-1}\left(\bigcap_{\alpha\in A}B_\alpha\right)=\bigcap_{\alpha\in A}f^{-1}(B_\alpha).
\]
\[
f(A\cap B)=f(A)\cap f(B)\text{ when }f\text{ is one-to-one.}
\]

\entry{Dense and separable}
\[
\text{Let }E\subseteq X.
\]
\[
E\text{ is dense in }X\Longleftrightarrow
\overline E=X
\Longleftrightarrow
\forall x\in X,\ \forall\text{ open }V\ni x,\ V\cap E\ne\varnothing.
\]
\[
X\text{ is separable if it has a countable dense subset.}
\]

\entry{Lindelof's Theorem}
\[
\text{Let }X\text{ be separable and }E\subseteq X.
\]
\[
\text{Every open cover }\{V_\alpha:\alpha\in A\}\text{ of }E
\text{ has a countable subcover.}
\]

\entry{Subspace topology}
\[
\text{Let }X\text{ be a metric space, }E\subseteq X,\ U,A\subseteq E.
\]
\[
U\text{ is open in }E
\Longleftrightarrow \exists\text{ an open set }V\text{ in }X\text{ such that }U=V\cap E.
\]
\[
A\text{ is closed in }E
\Longleftrightarrow \exists\text{ a closed set }B\text{ in }X\text{ such that }A=B\cap E.
\]

\newpage

\topic{Connectedness}

\entry{Disconnected and connected}
\[
\text{Let }(X,\rho)\text{ be a metric space, }E\subseteq X.
\]
\[
E\text{ is disconnected if }\exists\text{ nonempty disjoint sets }U,V\text{ open in }E
\text{ such that }E=U\cup V.
\]
\[
E\text{ is connected if it is not disconnected.}
\]
\[
\text{i.e. if }E=U\cup V,\ U,V\text{ disjoint and open in }E,\text{ then }U=\varnothing\text{ or }V=\varnothing.
\]

\entry{Open and closed in connected set}
\[
E\text{ is connected}\Longleftrightarrow
\varnothing\text{ and }E\text{ are the only subsets of }E\text{ that are both open and closed in }E.
\]

\entry{Intervals are connected}
\[
\text{In }(\R,d),\quad
E\text{ is connected}\Longleftrightarrow
E\text{ is an interval }(a,b),[a,b],(a,b],[a,b)
\]
\[
\text{or }E\text{ is a single point.}
\]

\entry{Continuity preserves connectedness}
\[
f\text{ continuous on connected }E
\Longrightarrow f(E)\text{ is connected.}
\]

\entry{Intermediate Value Theorem}
\[
f:E\to\R\text{ continuous on connected }E,\qquad f(x)<y<f(z)
\]
\[
\Longrightarrow \exists \xi\in E\text{ such that }f(\xi)=y.
\]

\entry{Path connected}
\[
E\text{ is path connected}
\Longleftrightarrow
\forall x,y\in E,\ \exists\text{ a continuous function }f:[0,1]\to E
\text{ such that }f(0)=x,\ f(1)=y.
\]

\entry{Path connected implies connected}
\[
E\text{ is path connected}\Longrightarrow E\text{ is connected.}
\]

\entry{Identity Theorem}
\[
\text{Let }f,g:I\to\R\text{ be analytic, where }I\text{ is an open interval.}
\]
\[
E=\{x\in I:f(x)=g(x)\}.
\]
\[
\text{If }E\text{ contains a cluster point of }E,\text{ then }f=g\text{ on }I.
\]

\newpage

\topic{Stone-Weierstrass}

\entry{Uniform metric space}
\[
\text{Let }X\text{ be compact and }C(X)=\{f:X\to\R:f\text{ continuous}\}.
\]
\[
\rho(f,g)=\norm{f-g}_\infty=\sup_{x\in X}|f(x)-g(x)|.
\]
\[
\text{Then }f_n\to f\text{ in }C(X)
\Longleftrightarrow f_n\to f\text{ uniformly on }X.
\]

\entry{Algebra}
\[
\text{Let }X\text{ be a set and }\mathcal A\text{ a set of functions from }X\text{ to }\R.
\]
\[
\mathcal A\text{ is an algebra if }f+g,\ fg,\ cf\in\mathcal A
\quad \forall f,g\in\mathcal A,\ \forall c\in\R.
\]
\[
\mathcal A\text{ separates points of }X
\Longleftrightarrow
\forall x,y\in X,\ x\ne y,\ \exists f\in\mathcal A\text{ such that }f(x)\ne f(y).
\]
\[
\mathcal A\text{ vanishes at no point of }X
\Longleftrightarrow
\forall x\in X,\ \exists f\in\mathcal A\text{ such that }f(x)\ne0.
\]
\[
\mathcal A\text{ is an algebra}\Longrightarrow\overline{\mathcal A}\text{ is an algebra.}
\]

\entry{Stone-Weierstrass Theorem}
\[
\text{Let }X\text{ be compact and }C(X)\text{ the uniform norm space.}
\]
\[
\mathcal A\subseteq C(X),\quad
\mathcal A\text{ is an algebra, separates points of }X,\text{ and vanishes at no point of }X
\]
\[
\Longrightarrow
\overline{\mathcal A}=C(X).
\]

\entry{Weierstrass Approximation Theorem}
\[
f\in C([a,b])
\Longrightarrow
\exists\text{ a sequence of polynomials }p_n\to f\text{ uniformly on }[a,b].
\]

\newpage

\topic{Partial Derivatives}

\entry{Partial Derivative}
\[
\frac{\partial f}{\partial x_j}(a)
=\lim_{h\to0}\frac{f(a+he_j)-f(a)}{h},
\qquad
D_jf(a)=\frac{\partial f}{\partial x_j}(a).
\]

\entry{$C^r$ functions}
\[
\text{Let }f:V\to\R^m,\quad V\subseteq\R^n,\quad r\in\N.
\]
\[
f\in C^r(V)\Longleftrightarrow
\text{all partial derivatives of }f\text{ of order }\le r
\text{ exist and are continuous on }V.
\]

\entry{Clairaut's Theorem}
\[
\text{Let }f:V\to\R,\quad V\text{ open in }\R^n.
\]
\[
f\in C^2(V)\Longrightarrow
\frac{\partial^2 f}{\partial x_j\partial x_k}(a)
=
\frac{\partial^2 f}{\partial x_k\partial x_j}(a).
\]

\entry{Interchange of limits and integrals}
\[
\text{Let }f_n:I\to\R\text{ be continuous and }f_n\to f\text{ uniformly on }I=[a,b].
\]
\[
\text{Then }f\text{ is continuous on }I,\qquad
\lim_{n\to\infty}\int_a^b f_n(x)\,dx
=\int_a^b f(x)\,dx.
\]

\entry{Interchange of derivative and integral}
\[
\text{Let }f:H\times\R\to\R,\quad H\times[a,b]\subseteq\R^2.
\]
\[
\text{If the relevant partial derivative is continuous, then }
\frac{d}{dy}\int_a^b f(x,y)\,dx
=\int_a^b \frac{\partial f}{\partial y}(x,y)\,dx.
\]

\topic{Linear Transformations}

\entry{Notation}
\[
A\in L(\R^n,\R^m),\qquad
A=[a_{ij}]_{m\times n}.
\]

\entry{Norm of matrix}
\[
\norm{A}=\sup_{\norm{x}=1}\norm{Ax}
=\sup_{x\ne0}\frac{\norm{Ax}}{\norm{x}}.
\]
\[
\norm{Ax}\le\norm{A}\norm{x},\qquad
\norm{AB}\le\norm{A}\norm{B},\qquad
\norm{A+B}\le\norm{A}+\norm{B}.
\]
\[
\rho(A,B)=\norm{A-B}\text{ is a metric on }L(\R^n,\R^m).
\]
\[
\norm{A}\le\left(\sum_{i,j}a_{ij}^2\right)^{1/2}.
\]

\entry{Cauchy-Schwarz}
\[
|u\cdot v|\le\norm{u}\norm{v}.
\]

\newpage

\topic{Differentiability}

\entry{Total derivative}
\[
\text{Let }V\subseteq\R^n\text{ be open, }a\in V.
\]
\[
f\text{ is differentiable at }a
\Longleftrightarrow
\exists T\in L(\R^n,\R^m)\text{ such that}
\]
\[
\lim_{h\to0}\frac{\norm{f(a+h)-f(a)-Th}}{\norm{h}}=0.
\]
\[
Df(a)=T.
\]

\entry{Continuous differentiability}
\[
\text{Let }V\subseteq\R^n\text{ be open and }f:V\to\R^m.
\]
\[
\text{If }f\text{ is continuously differentiable at }a,\text{ then }f\text{ is differentiable at }a.
\]
\[
\text{If }f\text{ is continuously differentiable on }V,\text{ then }f\in C^1(V).
\]

\entry{Continuously differentiable implies differentiable}
\[
f\text{ is continuously differentiable at }a
\Longrightarrow f\text{ is differentiable at }a.
\]

\entry{Chain rule}
\[
U\subseteq\R^n,\ V\subseteq\R^m\text{ open},\qquad
f:U\to\R^m,\quad g:V\to\R^k.
\]
\[
f(U)\subseteq V,\quad f\text{ differentiable at }a,\quad g\text{ differentiable at }f(a)
\]
\[
\Longrightarrow
g\circ f\text{ differentiable at }a,\qquad
D(g\circ f)(a)=Dg(f(a))Df(a).
\]

\topic{Mean Value Theorem}

\entry{Line segment}
\[
\text{Let }x,a\in\R^n.
\]
\[
\text{The line segment between }x\text{ and }a\text{ is }
L(x,a)=\{a+t(x-a):0\le t\le1\}.
\]

\entry{Convexity}
\[
E\subseteq\R^n\text{ is convex}
\Longleftrightarrow
L(x,y)\subseteq E,\quad \forall x,y\in E.
\]

\entry{Scalar function MVT}
\[
\text{Let }f:V\to\R\text{ be differentiable on }V,\quad V\subseteq\R^n\text{ open.}
\]
\[
V\text{ convex},\ x,a\in V
\Longrightarrow
\exists c\in L(x,a)\text{ such that}
\]
\[
f(x)-f(a)=Df(c)(x-a)=\nabla f(c)\cdot(x-a).
\]

\entry{Vector-valued MVT}
\[
\text{Let }f:V\to\R^m\text{ be differentiable on }V,\quad V\subseteq\R^n\text{ open.}
\]
\[
\forall x,a\in V\text{ with }L(x,a)\subseteq V,\quad
\exists c\in L(x,a)\text{ such that}
\]
\[
\norm{f(x)-f(a)}\le\norm{Df(c)}\norm{x-a}.
\]

\entry{Continuity of derivative map}
\[
\text{Let }V\subseteq\R^n\text{ be open and }f:V\to\R^m.
\]
\[
f\in C^1(V)
\Longleftrightarrow
Df:V\to L(\R^n,\R^m)\text{ is continuous.}
\]

\entry{Bound on compact and convex set}
\[
\text{Let }f:V\to\R^m\text{ be }C^1\text{ on }V,\quad V\subseteq\R^n\text{ open.}
\]
\[
\text{If }K\subseteq V\text{ is compact and convex, then }
\exists M\ge0\text{ such that}
\]
\[
\norm{f(x)-f(a)}\le M\norm{x-a},\qquad \forall x,a\in K.
\]

\newpage

\topic{Inverse Function Theorem and Implicit Function Theorem}

\entry{Invertible}
\[
\text{Let }B\in L(\R^n),\quad A\text{ is invertible.}
\]
\[
A+B\text{ and }B\in L(\R^n)\text{ with }\norm{B-A}\le\frac1{2\norm{A^{-1}}}
\Longrightarrow A+B\text{ is invertible and }\norm{(A+B)^{-1}}\le2\norm{A^{-1}}.
\]
\[
A\mapsto A^{-1}\text{ is a continuous map.}
\]

\entry{$C$-Lipschitz}
\[
f:X\to X\text{ is }C\text{-Lipschitz if }C<1
\text{ and }\rho(f(x),f(y))\le C\rho(x,y),\quad \forall x,y\in X.
\]

\entry{Contraction Principle}
\[
\text{Let }(X,\rho)\text{ be nonempty and complete metric space.}
\]
\[
\text{If }f:X\to X\text{ is }C\text{-Lipschitz for }C<1,
\text{ then it is a contraction and }\exists!x\in X\text{ such that }f(x)=x.
\]

\entry{Inverse Function Theorem}
\[
\text{Let }V\subseteq\R^n\text{ be open, }a\in V,\quad f:V\to\R^n
\text{ be }C^1\text{ on }V,\quad Df(a)\text{ invertible.}
\]
\[
\Longrightarrow
\exists U\ni a,\ W\ni f(a)\text{ open such that }f:U\to W\text{ is one-to-one and onto,}
\]
\[
f^{-1}:W\to U\text{ is }C^1\text{ and }
D(f^{-1})(f(x))=(Df(x))^{-1},\quad \forall x\in U.
\]

\entry{Partial Jacobian}
\[
\text{Let }F=(f_1,\ldots,f_m):U\to\R^m,\quad
U\subseteq\R^{n+m}\text{ open},\quad (x_0,y_0)\in U.
\]
\[
D_yF=
\left[\frac{\partial f_i}{\partial y_j}\right]_{m\times m},
\qquad
D_xF=
\left[\frac{\partial f_i}{\partial x_j}\right]_{m\times n}.
\]

\entry{Implicit Function Theorem}
\[
\text{Let }U\subseteq\R^{n+m}\text{ be open, }(a,b)\in U,\ c\in\R^m,
\quad f:U\to\R^m\text{ is }C^1\text{ on }U.
\]
\[
f(a,b)=c,\qquad D_yf(a,b)\text{ is invertible.}
\]
\[
\Longrightarrow
\exists V\ni a,\ W\ni b\text{ open and }g:V\to W\text{ such that for each }x\in V
\]
\[
\exists!y\in W\text{ with }(x,y)\in U\text{ and }f(x,y)=c,\qquad y=g(x).
\]
\[
g\in C^1(V),\qquad f(x,g(x))=c,\qquad
Dg(x)=-[D_yf(x,g(x))]^{-1}D_xf(x,g(x)).
\]

\newpage

\topic{Riemann-Lebesgue Theorem}

\entry{Upper and lower sums}
\[
\text{Let }\mathcal P\text{ be subrectangles of }P.
\]
\[
U(f,\mathcal P)=\sum_{M\in\mathcal P} M_f(M)v(M),
\qquad
L(f,\mathcal P)=\sum_{M\in\mathcal P} m_f(M)v(M),
\]
\[
M_f(M)=\sup_{x\in M}f(x),\qquad m_f(M)=\inf_{x\in M}f(x).
\]
\[
\int_A f=\sup_{\mathcal P}L(f,\mathcal P),\qquad
\int_A f=\inf_{\mathcal P}U(f,\mathcal P).
\]

\entry{Measure zero}
\[
E\text{ has measure zero in }\R^n
\Longleftrightarrow
\forall\varepsilon>0,\ \exists\text{ a countable cover }\{A_j\}_{j=1}^{\infty}
\text{ of }E\text{ by closed/open rectangles}
\]
\[
\text{such that }\sum_{j=1}^{\infty}v(A_j)<\varepsilon.
\]

\entry{Content zero}
\[
E\text{ has content zero if }
\forall\varepsilon>0,\ \exists\text{ a finite cover }\{A_1,\ldots,A_N\}
\text{ of }E
\]
\[
\text{by closed/open rectangles such that }\sum_{j=1}^{N}v(A_j)<\varepsilon.
\]

\entry{Riemann-Lebesgue Theorem}
\[
\text{Let }f:A\to\R\text{ be bounded, where }A\text{ is a closed rectangle in }\R^n.
\]
\[
f\text{ is integrable on }A
\Longleftrightarrow
\{x\in A:f\text{ is discontinuous at }x\}\text{ has measure zero.}
\]

\topic{Jordan Measurability and Fubini}

\entry{Characteristic function}
\[
\text{Let }A\text{ be a closed rectangle in }\R^n,\quad B\subseteq A.
\]
\[
\chi_B:A\to\R\text{ is the characteristic function,}\qquad
\chi_B(x)=
\begin{cases}
1,&x\in B,\\
0,&x\in A\setminus B.
\end{cases}
\]

\entry{Jordan-measurable}
\[
E\subseteq\R^n\text{ is Jordan-measurable if }E\text{ is bounded and }\partial E\text{ has measure }0.
\]

\entry{Integrability of }\(\chi_B\)
\[
\text{If }E\subseteq A\text{ is Jordan-measurable for a closed rectangle }A\subseteq\R^n,
\text{ then }\chi_E\text{ is integrable.}
\]
\[
\int_A\chi_B=\int_B1.
\]
\[
E\text{ has content zero}\Longleftrightarrow E\text{ is Jordan-measurable and }\int_A\chi_E=0.
\]
\[
E\text{ has measure zero and }E\text{ is Jordan-measurable}
\Longleftrightarrow
E\subseteq A\text{ for some closed rectangle }A\text{ and }\int_A\chi_E=0.
\]

\entry{Fubini's Theorem}
\[
\text{Let }A\subseteq\R^n,\ B\subseteq\R^m\text{ be closed rectangles,}
\quad f:A\times B\to\R\text{ be integrable on }A\times B.
\]
\[
\int_{A\times B}f
=\int_A\left(\int_B f(x,y)\,dy\right)dx
=\int_B\left(\int_A f(x,y)\,dx\right)dy.
\]
\[
\text{In particular, if for each }x\in A,\ \int_B f(x,y)\,dy\text{ exists,}
\text{ and for each }y\in B,\ \int_A f(x,y)\,dx\text{ exists,}
\]
\[
\text{then }\int_{A\times B}f
=\int_A\int_B f(x,y)\,dy\,dx
=\int_B\int_A f(x,y)\,dx\,dy.
\]

\end{document}
